% GENERATED FILE. Edit canonical lessons, then run npm run generate:books.
% canonical-source-hash: fnv1a64:8e293f4ffefd9e83
% canonical-lessons: GE-C506-torwart, GE-C506-stuermer, GE-C506-verteidiger, GE-C506-schiedsrichter, GE-C506-pfeifen

\chapter{Auf dem Platz}
\label{ch:ge-auf-dem-platz}
\hlchaptermodality{506}

\begin{chapteropening}
\textbf{By the end of this chapter:} \emph{I can talk about the players on a football pitch.}

Complete the last lesson of chapter 506: I can talk about the players on a football pitch.
\end{chapteropening}

\section[der Torwart]{der Torwart — goalkeeper}

\label{lesson:GE-C506-torwart}



\begin{quote}
Before the new one: say the German for to stretch, then the German for sore muscles.
\end{quote}

\subsection*{You'll want to know: der Torwart}
\textbf{der Torwart} — \textquotedblleft{}goalkeeper\textquotedblright{}.

\begin{quote}
Der Torwart hat den Ball sicher gefangen.
\end{quote}

\begin{quote}
\emph{The goalkeeper caught the ball safely.}
\end{quote}

Plural: \emph{die Torwarte}. Also \emph{der Torhüter}.

\begin{cousinweb}
\emph{Torwart} is built on \emph{Tor} \textquotedblleft{}gate, goal\textquotedblright{} (EtymDB) and \emph{Wart} \textquotedblleft{}keeper\textquotedblright{}, from \emph{warten}, inherited from Old High German \emph{wartēn} \textquotedblleft{}to watch, to look after\textquotedblright{} (EtymDB).
\end{cousinweb}

\begin{grammarlens}[title={What you've built}]
One more word for telling, explaining and describing at B1.
\end{grammarlens}

\subsection*{Your turn}
\begin{itemize}
\raggedright
  \item \emph{Say it:} \emph{der Torwart}
  \item \emph{Say it:} \emph{der Torwart}, once more
  \item \emph{Say it:} a sentence of your own with \emph{der Torwart}
  \item \emph{Recall:} say the German for to stretch, then the German for sore muscles, then say \emph{der Torwart} again
\end{itemize}

\begin{tcolorbox}[breakable,colback=teal!4,colframe=teal!35!black,title={Before you move on}]
What does der Torwart mean? (\textquotedblleft{}Goalkeeper\textquotedblright{}.) Say it once more.
\end{tcolorbox}
\section[der Stürmer]{der Stürmer — striker, forward}

\label{lesson:GE-C506-stuermer}



\begin{quote}
Before the new one: say the German for sore muscles, then the German for goalkeeper.
\end{quote}

\subsection*{You'll want to know: der Stürmer}
\textbf{der Stürmer} — \textquotedblleft{}striker, forward\textquotedblright{}.

\begin{quote}
Der Stürmer schoss zwei Tore.
\end{quote}

\begin{quote}
\emph{The striker scored two goals.}
\end{quote}

Plural: \emph{die Stürmer}.

\begin{cousinweb}
\emph{Stürmer} is built on \emph{stürmen} \textquotedblleft{}to storm\textquotedblright{}, from \emph{Sturm} \textquotedblleft{}storm, attack\textquotedblright{}, inherited from Old High German \emph{sturm}, from Proto-Germanic \emph{sturmaz} (EtymDB).
\end{cousinweb}

\begin{grammarlens}[title={What you've built}]
One more word for telling, explaining and describing at B1.
\end{grammarlens}

\subsection*{Your turn}
\begin{itemize}
\raggedright
  \item \emph{Say it:} \emph{der Stürmer}
  \item \emph{Say it:} \emph{der Stürmer}, once more
  \item \emph{Say it:} a sentence of your own with \emph{der Stürmer}
  \item \emph{Recall:} say the German for sore muscles, then the German for goalkeeper, then say \emph{der Stürmer} again
\end{itemize}

\begin{tcolorbox}[breakable,colback=teal!4,colframe=teal!35!black,title={Before you move on}]
What does der Stürmer mean? (\textquotedblleft{}Striker, forward\textquotedblright{}.) Say it once more.
\end{tcolorbox}
\section[der Verteidiger]{der Verteidiger — defender}

\label{lesson:GE-C506-verteidiger}



\begin{quote}
Before the new one: say the German for goalkeeper, then the German for striker.
\end{quote}

\subsection*{You'll want to know: der Verteidiger}
\textbf{der Verteidiger} — \textquotedblleft{}defender\textquotedblright{}.

\begin{quote}
Der Verteidiger stoppte den Stürmer kurz vor dem Tor.
\end{quote}

\begin{quote}
\emph{The defender stopped the striker shortly before the goal.}
\end{quote}

Plural: \emph{die Verteidiger}. Also a defence lawyer.

\begin{cousinweb}
\emph{Verteidiger} is built on \emph{verteidigen} \textquotedblleft{}to defend\textquotedblright{}, usually traced to Middle High German \emph{vertagedingen} \textquotedblleft{}to plead before a court\textquotedblright{}; EtymDB lists the verb without a source.
\end{cousinweb}

\begin{grammarlens}[title={What you've built}]
One more word for telling, explaining and describing at B1.
\end{grammarlens}

\subsection*{Your turn}
\begin{itemize}
\raggedright
  \item \emph{Say it:} \emph{der Verteidiger}
  \item \emph{Say it:} \emph{der Verteidiger}, once more
  \item \emph{Say it:} a sentence of your own with \emph{der Verteidiger}
  \item \emph{Recall:} say the German for goalkeeper, then the German for striker, then say \emph{der Verteidiger} again
\end{itemize}

\begin{tcolorbox}[breakable,colback=teal!4,colframe=teal!35!black,title={Before you move on}]
What does der Verteidiger mean? (\textquotedblleft{}Defender\textquotedblright{}.) Say it once more.
\end{tcolorbox}
\section[der Schiedsrichter]{der Schiedsrichter — referee}

\label{lesson:GE-C506-schiedsrichter}



\begin{quote}
Before the new one: say the German for striker, then the German for defender.
\end{quote}

\subsection*{You'll want to know: der Schiedsrichter}
\textbf{der Schiedsrichter} — \textquotedblleft{}referee\textquotedblright{}.

\begin{quote}
Der Schiedsrichter zeigte dem Spieler die gelbe Karte.
\end{quote}

\begin{quote}
\emph{The referee gave the player a yellow card.}
\end{quote}

Plural: \emph{die Schiedsrichter}. In speech also \emph{der Schiri}.

\begin{cousinweb}
\emph{Schiedsrichter} is built on \emph{Schied} \textquotedblleft{}decision\textquotedblright{}, from \emph{scheiden} \textquotedblleft{}to separate, to decide\textquotedblright{}, and \emph{Richter} \textquotedblleft{}judge\textquotedblright{} (EtymDB lists the German words).
\end{cousinweb}

\begin{grammarlens}[title={What you've built}]
One more word for telling, explaining and describing at B1.
\end{grammarlens}

\subsection*{Your turn}
\begin{itemize}
\raggedright
  \item \emph{Say it:} \emph{der Schiedsrichter}
  \item \emph{Say it:} \emph{der Schiedsrichter}, once more
  \item \emph{Say it:} a sentence of your own with \emph{der Schiedsrichter}
  \item \emph{Recall:} say the German for striker, then the German for defender, then say \emph{der Schiedsrichter} again
\end{itemize}

\begin{tcolorbox}[breakable,colback=teal!4,colframe=teal!35!black,title={Before you move on}]
What does der Schiedsrichter mean? (\textquotedblleft{}Referee\textquotedblright{}.) Say it once more.
\end{tcolorbox}
\section[pfeifen]{pfeifen — to whistle}

\label{lesson:GE-C506-pfeifen}



\begin{quote}
Before the new one: say the German for defender, then the German for referee.
\end{quote}

\subsection*{You'll want to know: pfeifen}
\textbf{pfeifen} — \textquotedblleft{}to whistle\textquotedblright{}.

\begin{quote}
Der Schiedsrichter pfeift, und das Spiel beginnt.
\end{quote}

\begin{quote}
\emph{The referee whistles, and the game begins.}
\end{quote}

Irregular: \emph{pfiff, hat gepfiffen}. \emph{Anpfeifen}, \emph{abpfeifen} — to start, to end a match.

\begin{cousinweb}
\emph{pfeifen} is inherited from Middle High German \emph{pfīfen} and Old High German \emph{phīfōn}, an early borrowing from Latin \emph{pīpāre} \textquotedblleft{}to peep, to chirp\textquotedblright{} (EtymDB).
\end{cousinweb}

\begin{grammarlens}[title={What you've built}]
One more word for telling, explaining and describing at B1.
\end{grammarlens}

\subsection*{Your turn}
\begin{itemize}
\raggedright
  \item \emph{Say it:} \emph{pfeifen}
  \item \emph{Say it:} \emph{pfeifen}, once more
  \item \emph{Say it:} a sentence of your own with \emph{pfeifen}
  \item \emph{Recall:} say the German for defender, then the German for referee, then say \emph{pfeifen} again
\end{itemize}

\begin{tcolorbox}[breakable,colback=teal!4,colframe=teal!35!black,title={Before you move on}]
What does pfeifen mean? (\textquotedblleft{}To whistle\textquotedblright{}.) Say it once more.
\end{tcolorbox}
